\chapter{Introduction}

	\noindent
	The general proccess, begins KNOWING perfectly the set we are talking about. We need to understand HOW their elements are constructed, and their combinatorics.
	\\\\
	
	\noindent
	In Theory of Languages, formal Languages are two sets: one set contains the alphabet of symbols we are going to use; and the other set contains all chains we decided the language contains, using JUST the symbols of the previous alphabet to create the chains. We are going to turn tables. NOW, sets are languages. We are going to define a particular alphabet ( finite or infinite ) of symbols for our set, and the elements of the set, are going to be CHAINS of those symbols... finite chains, or INFINITE chains. We are going to call our symbols, LAMBDAS, and each element is going to be in the format of a t-upla of lambdas:\\\\
	\textless ** \textit{See unnecesary comment 0001, in case you want to know WHY we call them lambdas} \textgreater\\
	$$ (\lambda_{1}, \lambda_{2}, \lambda_{3}, \lambda_{4}, ... , \lambda_{m}) $$\\
	$ \lambda_{1} $ and $\lambda_{4}$ can be the same symbol of our alphabet. Lambdas are more like variables that can contain any possible symbol of our alphabet.
	\\\\
	
	\noindent
	The second step is designing a Construction LJA for our set. And this is hard. We need to design the alphabet, and the properties of the chains, to fit our set inside a particular CLJA. Or another set, that contains our studied set. We want just to prove it HAS NOT a quantity of elements bigger than $\mathbb{N}$ has. So a superset of it works well for us too.
	\\\\
	
	\noindent
	Once we had designed our Language, the set of representations, for our set, that could be inside a CLJA, it will define ANOTHER LANGUAGE, called LCF (wait for it, it is explained after, in this document). LCF only has chains of finite size, and thanks we have it produced by a CLJA, the bijection with N is automatic. Direct and inverse funtions. Well defined mathemathically.
	\\\\
	
	\noindent
	The system of relations is created between our set of representations (the first language), and the language LCF defined by the CLJA (our second language).
	\\\\
	
	\noindent
	We need to provide any possible tool, to prove our studied set, HAS NOT A CARDINAL bigger than our first language (its set of chains). And a second tool to prove our LCF language has not a cardinal bigger than $\mathbb{N}$. It is provided by the bijection the CLJA creates with N, called: FLJA... Function LJA.
	\\\\
	
	\noindent
	CLJAs have an important rule: the alphabet can be infinite, but we need to provide a bijection with $mathbb{N}$. We can deal with sets with any possible $\aleph$ as its cardinal, but the alphabet we use must have cardinal $\aleph_{0}$, and we need to define the bijection between $\mathbb{N}$ and the alphabet we use, for the language we use, to represent our studied set. Inside FLJAs there is a subfunction called ``Translation function'', and it uses THAT bijection, inside it.
	\\\\
	
	\noindent
	We are not talking anymore about CLJAs. And it hurts me a lot. Why? It could take several days, or hundreds of pages. I will prove ONLY that the bijection between our particular case of LCF and N, MUST EXISTS, based in a very well known result. And we must accept the deal, you will NOT try to change anything. You will accept our representation language, and LCF as they are. Even if they are a little weird. The exaplanation is based in the CLJA, called CLJA-FTC, that we designed for $\mathcal{P}(\mathbb{N})$ vs $\mathbb{N}$. And we don't need the CLJA to explain the future properties of both languages and what they can do.
	\\\\
	
	\noindent
	\textless ** \textit{See unnecesary comment 0002, to read curious facts about CLJAs, under your own risk :D} \textgreater
	\\\\
	
	\noindent
	The general strategy is:\\
	$$ Card(\{ \text{ORIGINAL SET} \}) \ngtr Card(\{ \text{REPRESENTATION LANGUAGE} \}) \ngtr Card(LCF) \ngtr Card(\mathbb{N}) $$
	And show it by any possible way.
	\\\\
	
	\noindent
	In our particular case:
	
	
	
		
	
	
	


\newpage